\documentclass[11pt]{article}
\usepackage[a4paper, top=18mm, bottom=18mm, left=20mm, right=20mm]{geometry}
\usepackage[T2A]{fontenc}
\usepackage[utf8]{inputenc}
\usepackage[ukrainian]{babel}
\usepackage{graphicx}
\usepackage[export]{adjustbox}
\usepackage{amsmath}
\usepackage{amssymb}
\graphicspath{{/app/Quiz/Scripts/2023/ProgAll/15BinTree/}{/app/Quiz/Scripts/2021/Mod2021_3/BinTree/}{/app/Quiz/Scripts/2020/Mod2020_4/BinTree/}}
\setlength{\parindent}{0pt}
\setlength{\parskip}{6pt}
\pagestyle{empty}
\begin{document}
%begin
{\large\textbf{Контрольна робота}}

\textbf{Варіант \No\,10}\hfill \textit{ПІБ:}\,\underline{\hspace{60mm}}
\vspace{4mm}

%beginitem
1. Що буде виведено за виконання фрагменту коду?

%code
#include <iostream>

int h(int a, int b){
    int c = 69;
    if (b >= -3)
        return 3;
    if (a <= -4)
         c = 1;
    else 
        return 5;
    return c;
}

int main(){
    std::cout << h(9, 9);
    return 0;
}
%code
\vspace{5mm}
%enditem

%beginitem
2. Що буде виведено за виконання фрагменту коду?
try:
    res = "3">=3.0
    if isinstance(res, bool):
        print(f'bool;{res}')
    elif isinstance(res, int):
        print(f'int;{res}')
    elif isinstance(res, float):
        print(f'float;{res:.2f}')
    else:
        print('???')
except:
    print('Z')
\vspace{5mm}
%enditem

%beginitem
3. %goodpath:Scripts/2018/Mod2018_1/03-good.txt
Відмітити все, що є коректним оператором С++:
( 1 ) cout<<(3=!5);

( 2 ) bool f(int a);

( 3 ) cin<<x;

( 4 ) j=i;
\vspace{5mm}
%enditem

%beginitem
4. 

\begin{center}
	\adjustimage{max size={0.8\linewidth}{0.8\paperheight}}
	{../BinTree/trees_exp/59.png}
\end{center}
Указати послідовність міток вершин дерева, зображеного на рис., що утвориться в результаті обходу дерева в прямому порядку. Мітки записувати через пробіл.\par Обхід дерева починається з кореня (на рис. це верхня вершина).
%answer:59 прямому
\vspace{5mm}
%enditem

%beginitem
5. Що буде виведено за виконання фрагменту коду?
def f(n):
    print('f', end='')
    return n

try:
    print(f(9) and f(-8))
except:
    print('ERR')
\vspace{5mm}
%enditem

%beginitem
6. Що буде виведено за виконання фрагменту коду?
d=[10,11,12,13,14]; del d[-6:-6]; print(d)\vspace{5mm}
%enditem

%beginitem
7. Що буде виведено за виконання фрагменту коду?

%code
print(5/5%5*5)
%code
\vspace{5mm}
%enditem

%beginitem
8. Що буде виведено за виконання фрагменту коду?
try:
    print(3, end=';')
    print(1j>=3j, end=';')
    print(5, end=';')
except TypeError:
    print(0, end=';')
except ValueError:
    print(9, end=';')
except:
    print('Z', end=';')
else:
    print(6, end=';')
finally:
    print(4)
\vspace{5mm}
%enditem

%beginitem
9. %goodpath:Scripts/2019/Mod2019_1/Lex/03-good.txt
Відмітити все, що є коректним оператором С++:
( 1 ) if a<b: a=0 else: b=0

( 2 ) x=3**1

( 3 ) if ('x'><'y'): a=6

( 4 ) --n
\vspace{5mm}
%enditem

%beginitem
10. %goodpath:Scripts/2018/Mod2018_1/01-good.txt
Відмітити все, що є однією правильною лексемою:
( 1 ) 0,5

( 2 ) LIFO

( 3 ) x+23

( 4 ) x23
\vspace{5mm}
%enditem

%end
\newpage
\end{document}

